% =====================================================================
%  PACE : unified equation set (AAAI-2027)
%  Single shared notation across problem setup, diffusion policy,
%  all interactive-IL baselines, and PACE.
%  ASCII-only math. Every block is preceded by a % intent comment.
%  PACE is presented as a strict generalization of the query-based
%  baselines: baselines answer "WHEN to query"; PACE answers
%  "WHEN + WHICH failures + WHERE to place the corrective demo".
% =====================================================================

% ---------------------------------------------------------------------
%  NOTATION TABLE  (symbol -> meaning). Keep identical everywhere.
% ---------------------------------------------------------------------
% -- Environment / policy --
%   s in S            state (privileged proprio for state-mode; image obs for image-mode)
%   a in A            action (discrete move in T1; rel_joint_pos delta in T2-T5)
%   pi_theta          novice (learner) policy, parameters theta
%   pi_star           expert oracle (A*/BFS in T1; PPO/motion-planner in T2-T5)
%   P(s'|s,a), R, H   transition kernel, reward, horizon
%   gamma_env         MDP discount (distinct from the memory discount gamma below)
%   tau_traj          a trajectory (s_0,a_0,...); rho_pi its state-visitation dist.
% -- Datasets / loop --
%   D                 aggregated demonstration dataset (state-action pairs)
%   r                 interactive-IL round index (one prescribed demo per round)
%   n_d               retrain cadence (retrain from scratch every n_d demos)
%   B                 expert-query budget (max added demos); q = #queries used
% -- Diffusion policy --
%   x_0 = a           clean action (target of denoising), a in R^{da}
%   k in {1..K_dif}   diffusion timestep; K_dif = num_train_timesteps
%   eps ~ N(0,I)      injected noise; alpha_bar_k the DDPM cumulative schedule
%   x_k               noised action at diffusion step k
%   v_k               v-prediction target
%   eps_theta,v_theta network prediction (noise / v param.) conditioned on obs
%   L_dif(s,a)        per-pair diffusion (denoising) loss  == UNCERTAINTY SIGNAL
%   ell_t = L_dif(s_t, a^{nov}_t)   diffusion loss at rollout step t
%   F_hat             empirical CDF of training-set diffusion losses
%   alpha             CDF quantile level for the Diff-DAgger threshold (0.99)
% -- Interactive-IL query machinery (shared by ALL baselines) --
%   a^{nov}_t         novice's first executed action at step t
%   a^{exp}_t         expert action at s_t (both in comparable rel_joint_pos space)
%   Query(s_t)        binary predicate: 1 => expert intervenes at t*
%   t* = min{t: Query(s_t)=1}   first intervention step of the round
%   tau_sd, tau,...   per-method thresholds (subscripts name the method)
% -- MC-dropout / ensemble --
%   N_mc              # stochastic samples (MC-dropout / diffusion draws)
%   a^{(i)}_t         i-th stochastic action sample at s_t
%   M                 # ensemble members; a^{[m]}_t member m's action
%   p_hat_t           in-ball fraction (agreement with expert)
% -- Success-Q (ThriftyDAgger) --
%   Q_psi(s,a)        goal-conditioned success value, params psi
%   y_t               Monte-Carlo success-to-go BCE target
%   risk(s,a)=1-Q     task-risk term
%   delta_h, beta_h   novelty / risk quantile thresholds (level 1-alpha_h)
% -- PACE PARTITION/PRIORITIZE (per round, on the round's failures) --
%   F_r={f_1..f_N}    the round's failure set; i,j index failures; N=|F_r|
%   t*_i,T_i          peak-loss step / total steps of failure i
%   peak_i,mean_i     max / mean diffusion loss over failure i
%   p^tee_i,theta_i   object (tee) planar pose + yaw of failure i
%   p^tcp_i,q^full_i  end-effector pos / full robot qpos at t*_i
%   rho_i,delta_i     progress t*_i/T_i ; contact distance ||p^tcp_i-p^tee_i||
%   phi_i in R^6      geometric failure descriptor  (featurizer phi(.))
%   v_i, v_tilde_i    raw R3M frame embedding / PCA-reduced visual embedding
%   psi_i             clustering feature = v_tilde_i (visual) else phi_i
%   X, X_tilde        stacked features / z-standardized features
%   C_1..C_m          clusters; k* silhouette-optimal cluster count
%   rep(C)            cluster representative (member nearest centroid)
%   C_star,C_tgt      dominant cluster / memory-rotated target cluster
%   S                 diversity-selected failure subset, |S|<=kappa
%   kappa             selection cap == top_k (default 3)
%   Pmem(c)           cross-round recency penalty at pose c
%   gamma,sigma_mem   memory discount / bandwidth ; lambda memory weight
% -- PACE PRESCRIBE --
%   A                 anchor = rep(C_tgt) geometry
%   cmd               LLM SceneCommand (target tee pose, tcp, label)
%   g(.)              prescribe map: SceneCommand -> ResetSpec (reset dist.)
%   xi                corrective reset / sub-task-entry spec (start of demo)
%   clamp_W(.)        projection onto KAG workspace bounds W
% -- Metrics --
%   SR               held-out success rate; SR_target = 0.90
%   q(SR_target)     #expert queries to reach target SR
%   Eff             demo-efficiency = AUC of SR vs #demos
%   Cov             coverage of collected demos over the task space
% ---------------------------------------------------------------------


% =====================================================================
%  PART I -- PROBLEM SETUP
% =====================================================================

% (1) Sequential decision process (MDP; POMDP if s is a partial image obs).
\begin{equation}
\mathcal{M} = \big(\, \mathcal{S},\ \mathcal{A},\ P(s'\mid s,a),\ R(s,a),\ H,\ \gamma_{\mathrm{env}} \,\big),
\qquad
\pi_\theta:\ \mathcal{S}\to\Delta(\mathcal{A}),
\label{eq:mdp}
\end{equation}

% (2) Expert oracle and the demonstration dataset it induces (roll pi_star).
\begin{equation}
\pi^\star=\arg\max_{\pi}\ \mathbb{E}_{\tau\sim\pi}\!\Big[\textstyle\sum_{t=0}^{H}\gamma_{\mathrm{env}}^{\,t} R(s_t,a_t)\Big],
\qquad
\mathcal{D}=\big\{(s_t,a_t)\ :\ a_t=\pi^\star(s_t)\big\}.
\label{eq:expert}
\end{equation}

% (3) Behavior-cloning objective: match the expert on the aggregated data.
\begin{equation}
\theta^\star=\arg\min_{\theta}\ \mathbb{E}_{(s,a)\sim\mathcal{D}}\big[\,\mathcal{L}_{\mathrm{BC}}\big(\pi_\theta(\cdot\mid s),\,a\big)\,\big],
\qquad
\mathcal{L}_{\mathrm{BC}}=-\log\pi_\theta(a\mid s)\ \ \text{(disc.)}\ \ \text{or}\ \ \|\pi_\theta(s)-a\|_2^2\ \ \text{(cont.)}.
\label{eq:bc}
\end{equation}

% (4) Interactive-IL query loop: DAgger-style aggregation with a query gate.
%     A round rolls pi_theta, flags an intervention point t*, the expert takes
%     over there, and the resulting expert segment is the single new demo.
\begin{align}
\text{roll } \pi_{\theta_r}\ &\Rightarrow\ t^\star=\min\{t:\ \mathrm{Query}(s_t)=1\}, \nonumber\\
\mathcal{D}_{r+1}&=\mathcal{D}_r\ \cup\ \big\{(s_t,\pi^\star(s_t)):\ t\ge t^\star\big\},\qquad
\theta_{r+1}=\arg\min_{\theta}\ \mathbb{E}_{\mathcal{D}_{r+1}}[\mathcal{L}_{\mathrm{BC}}],
\label{eq:iil-loop}
\end{align}
% one demo per round; retrain from scratch every n_d demos; stop at SR_target or budget B.


% =====================================================================
%  PART II -- DIFFUSION POLICY + ITS PER-STEP UNCERTAINTY SIGNAL
% =====================================================================

% (5) Forward noising of the action target a=x_0 over K_dif diffusion steps.
\begin{equation}
x_k=\sqrt{\bar\alpha_k}\,x_0+\sqrt{1-\bar\alpha_k}\,\epsilon,
\qquad \epsilon\sim\mathcal{N}(0,I),\quad k\in\{1,\dots,K_{\mathrm{dif}}\}.
\label{eq:forward}
\end{equation}

% (6) Training objective: v-prediction denoising (equiv. DDPM eps-prediction).
%     v_k is the DDPM v-target; conditioning on obs s is suppressed for brevity.
\begin{equation}
v_k=\sqrt{\bar\alpha_k}\,\epsilon-\sqrt{1-\bar\alpha_k}\,x_0,
\qquad
\mathcal{L}_{\mathrm{DP}}(\theta)=\mathbb{E}_{k,\epsilon,(s,a)}\big[\big\|v_\theta(x_k,k,s)-v_k\big\|_2^2\big].
\label{eq:vpred}
\end{equation}

% (7) The per-pair diffusion loss reused as the UNCERTAINTY SIGNAL.
%     Averaged over noise draws (and, optionally, all k); ell_t is its value
%     at the novice's own executed action along a rollout.
\begin{equation}
L_{\mathrm{dif}}(s,a)=\mathbb{E}_{k,\epsilon}\big[\big\|v_\theta(x_k,k,s)-v_k\big\|_2^2\big],
\qquad
\ell_t \;=\; L_{\mathrm{dif}}\big(s_t,\,a^{\mathrm{nov}}_t\big).
\label{eq:diffsignal}
\end{equation}


% =====================================================================
%  PART III -- BASELINE QUERY PREDICATES (a SINGLE decision rule family)
%  Every baseline is  Query(s_t)=1[ score_method(s_t) crosses threshold ].
%  Actions compared in the shared rel_joint_pos space so ||.|| is consistent.
% =====================================================================

% (8) SafeDAgger*: query when novice/expert action discrepancy exceeds tau_sd.
%     (Robot: L2 in joint-delta space; T1: fraction of off-optimal-mask steps.)
\begin{equation}
\mathrm{Query}_{\mathrm{Safe}}(s_t)=\mathbf{1}\big[\ \|a^{\mathrm{nov}}_t-a^{\mathrm{exp}}_t\|_2 > \tau_{\mathrm{sd}}\ \big],
\qquad
d^{\mathrm{maze}}_{\mathrm{safe}}=\tfrac1N\!\sum_{t}\mathbf{1}\!\big[\,\mathrm{opt}_t[a^{\mathrm{nov}}_t]<\tfrac12\,\big].
\label{eq:safe}
\end{equation}

% (9) DropoutDAgger: N_mc stochastic samples; query when the in-tau-ball
%     agreement fraction p_hat_t with the expert drops below p.
\begin{align}
\hat p_t=\tfrac{1}{N_{\mathrm{mc}}}\sum_{i=1}^{N_{\mathrm{mc}}}\mathbf{1}\big[\ \|a^{(i)}_t-a^{\mathrm{exp}}_t\|_2\le\tau\ \big],
\qquad
\mathrm{Query}_{\mathrm{Drop}}(s_t)=\mathbf{1}\big[\ \hat p_t<p\ \big].
\label{eq:dropout}
\end{align}

% (10) EnsembleDAgger: query on ensemble DOUBT (variance) OR mean DISCREPANCY.
\begin{align}
\mathrm{doubt}_t&=\tfrac1{d_a}\!\sum_{j=1}^{d_a}\mathrm{Var}_{m}\big(a^{[m]}_{t,j}\big),
&
\mathrm{disc}_t&=\tfrac1M\!\sum_{m=1}^{M}\big\|a^{[m]}_t-a^{\mathrm{exp}}_t\big\|_2, \nonumber\\
\mathrm{Query}_{\mathrm{Ens}}(s_t)&=\mathbf{1}\big[\ \mathrm{disc}_t>\tau\ \ \lor\ \ \mathrm{doubt}_t>\chi\ \big]. &&
\label{eq:ensemble}
\end{align}

% (11) ThriftyDAgger: query on NOVELTY (ensemble doubt) OR RISK=1-Q.
%      Success-Q trained by BCE on a Monte-Carlo success-to-go target y_t;
%      thresholds are budget-alpha_h auto-calibrated (1-alpha_h)-quantiles.
\begin{align}
y_t&=\mathbf{1}[\mathrm{success}]\cdot\gamma^{\,n-1-t},
\qquad
\mathcal{L}_Q=\mathrm{BCE}\big(\sigma(Q_\psi(s_t,a_t)),\,y_t\big),
\qquad
\mathrm{risk}(s,a)=1-\sigma\!\big(Q_\psi(s,a)\big), \nonumber\\
\delta_h&=\mathrm{Quantile}_{1-\alpha_h}\{\mathrm{doubt}\},\quad
\beta_h=\mathrm{Quantile}_{1-\alpha_h}\{\mathrm{risk}\},\quad
\mathrm{Query}_{\mathrm{Thr}}(s_t)=\mathbf{1}\big[\ \mathrm{doubt}_t>\delta_h\ \lor\ \mathrm{risk}_t>\beta_h\ \big].
\label{eq:thrifty}
\end{align}

% (12) Diff-DAgger: native diffusion-loss CDF quantile held for K windowed steps.
%      Threshold = alpha-quantile of the training-loss CDF, recalibrated each retrain.
\begin{equation}
\eta_\alpha=\hat F^{-1}(\alpha),
\qquad
\mathrm{Query}_{\mathrm{Diff}}(s_t)=\mathbf{1}\!\left[\ \sum_{u=t-K+1}^{t}\mathbf{1}\big[\ \ell_u>\eta_\alpha\ \big]\ \ge\ K\ \right]
\;=\;\mathbf{1}\big[\ \hat F(\ell_u)>\alpha\ \ \forall u\in(t-K,t]\ \big].
\label{eq:diffdagger}
\end{equation}

% (13) Uniform-random control ("Stagger"): query at a pre-drawn random step.
\begin{equation}
t_{\mathrm{stag}}\sim\mathrm{Unif}\{0,\dots,H-1\},
\qquad
\mathrm{Query}_{\mathrm{Rand}}(s_t)=\mathbf{1}\big[\ t\ge t_{\mathrm{stag}}\ \big].
\label{eq:stagger}
\end{equation}


% =====================================================================
%  PART IV -- PACE  (Perceive -> Assess -> Choose -> Execute)
%  Generalization: baselines pick WHEN (a single t*); PACE additionally
%  partitions the round's failures, selects a diverse WHICH-set, and
%  prescribes WHERE the corrective demo begins.
% =====================================================================

% (14) PERCEIVE. Failure set of the round; each failure keeps its peak-loss
%      step t*_i = argmax_t ell_t (the moment of maximal diffusion uncertainty),
%      which is exactly the Diff-DAgger signal (12) reused as the analysis frame.
\begin{equation}
\mathcal{F}_r=\{f_1,\dots,f_N\},
\qquad
t^\star_i=\arg\max_t \ell^{\,(i)}_t,\quad
\mathrm{peak}_i=\max_t \ell^{\,(i)}_t,\quad
\rho_i=\frac{t^\star_i}{\max(1,T_i)},\quad
\delta_i=\big\|p^{\mathrm{tcp}}_i-p^{\mathrm{tee}}_i\big\|_2.
\label{eq:perceive}
\end{equation}

% (15) Featurizer phi(.): 6-D geometric descriptor (quaternion-free);
%      optional visual embedding v_i (R3M) reduced by PCA to v_tilde_i.
%      psi_i is the clustering feature actually used; X_tilde its z-standardization.
\begin{align}
\phi_i&=\big[\,p^{\mathrm{tee}}_{i,x},\ p^{\mathrm{tee}}_{i,y},\ \sin\theta_i,\ \cos\theta_i,\ \rho_i,\ \delta_i\,\big]\in\mathbb{R}^6,
\qquad
\tilde v_i=\mathrm{PCA}_{k'}\big(v_i\big),\ \ k'=\min(16,N-1,d), \nonumber\\
\psi_i&=\begin{cases}\tilde v_i & \text{visual mode}\\[2pt]\phi_i & \text{geometric (default)}\end{cases},
\qquad
X=[\psi_1;\dots;\psi_N],\qquad
\tilde X_i=\big(\psi_i-\mu\big)\oslash\sigma,\quad \sigma_j\!\leftarrow\!\max(\sigma_j,10^{-8}).
\label{eq:phi}
\end{align}

% (16) PARTITION. Cluster the standardized features; choose k by max silhouette.
%      Per-cluster geometry (centroid, circular-mean yaw, mean peak-loss) is
%      always computed on the RAW tee pose; the representative is the member
%      nearest the cluster mean in feature space.
\begin{align}
k^\star&=\arg\max_{k\in[2,\,k_{\max}]}\ \mathrm{sil}(k),\qquad k_{\max}=\max\!\big(2,\min(6,N-1)\big),\nonumber\\
c^{xy}_{C}&=\tfrac{1}{|C|}\!\sum_{i\in C}p^{\mathrm{tee}}_{i,xy},\quad
\bar\theta_{C}=\operatorname{atan2}\!\Big(\!\sum_{i\in C}\!\sin\theta_i,\ \sum_{i\in C}\!\cos\theta_i\Big),\quad
\bar L_{C}=\tfrac{1}{|C|}\!\sum_{i\in C}\mathrm{peak}_i,\nonumber\\
\mathrm{rep}(C)&=\arg\min_{i\in C}\ \big\|\tilde X_i-\bar X_{C}\big\|_2,\qquad
\bar X_{C}=\tfrac1{|C|}\!\sum_{i\in C}\tilde X_i.
\label{eq:partition}
\end{align}

% (17) Dominant cluster: lexicographic on (size, mean peak-loss, -min episode id).
\begin{equation}
C^\star=\operatorname*{arg\,lexmax}_{C\in\{C_1..C_m\}}\ \big(\ |C|,\ \bar L_{C},\ -\!\min_{i\in C}\mathrm{eid}_i\ \big).
\label{eq:dominant}
\end{equation}

% (18) Cross-round MEMORY: Gaussian, recency-discounted coverage penalty, and
%      the memory-rotated TARGET cluster (among clusters within one member of
%      dominant). Rotation avoids re-prescribing an already-covered mode.
\begin{align}
P_{\mathrm{mem}}(c)&=\sum_{i:\,(r_i,c_i)\in\mathrm{Mem}}\gamma^{\,\max(0,\,r-r_i)}\exp\!\Big(-\tfrac{\|c-c_i\|_2^2}{2\sigma_{\mathrm{mem}}^2}\Big),\nonumber\\
C_{\mathrm{tgt}}&=\arg\max_{C:\,|C|\ge|C^\star|-1}\ \Big(\ \bar L_{C}\ -\ \lambda\,P_{\mathrm{mem}}\big(c^{xy}_{C}\big)\ \Big),\qquad \gamma=0.6,\ \sigma_{\mathrm{mem}}=0.06,\ \lambda=1.
\label{eq:memory}
\end{align}

% (19) PRIORITIZE. Diversity selection = k-center / farthest-point (max-min)
%      with the target representative FORCED in and the worst-loss seed added.
%      This is the coverage/diversity objective PACE adds over WHEN-only rules.
\begin{align}
S_0&=\big\{\,\mathrm{rep}(C_{\mathrm{tgt}})\,\big\}\ \cup\ \big\{\,\arg\max_i \mathrm{peak}_i\,\big\}, \nonumber\\
S&\leftarrow S\ \cup\ \Big\{\ \arg\max_{i\notin S}\ \min_{j\in S}\ \big\|\tilde X_i-\tilde X_j\big\|_2\ \Big\}\quad\text{until}\ |S|=\kappa,\qquad \kappa=3, \nonumber\\
S^\star&=\arg\max_{S\subseteq\mathcal{F}_r,\ |S|\le\kappa}\ \min_{\substack{i,j\in S\\ i\ne j}}\ \big\|\tilde X_i-\tilde X_j\big\|_2\quad\text{s.t.}\ \ \mathrm{rep}(C_{\mathrm{tgt}})\in S.
\label{eq:diversity}
\end{align}
% Eq. (19) top two lines = the greedy realization; last line = the max-min
% coverage objective it approximates (forced-representative constraint).

% (20) PRESCRIBE. The anchor is the target representative; the VLM/LLM emits a
%      SceneCommand cmd; the prescribe map g(.) turns it into a corrective reset
%      / sub-task-entry distribution xi (the START STATE of the one new demo),
%      projected onto the KAG workspace bounds W. A(nchor)-relative perturbation
%      is L2-capped in position and clamped in yaw.
\begin{align}
A&=\big(p^{\mathrm{tee}}_A,\ \theta_A,\ q^{\mathrm{full}}_A,\ p^{\mathrm{tcp}}_A,\ t^\star_A\big)=\mathrm{rep}(C_{\mathrm{tgt}}),
\qquad
\mathrm{cmd}=\mathrm{VLM/LLM}\big(A,\ \text{frames}(t^\star_A),\ \mathrm{KAG}\big),\nonumber\\
\Delta_{xy}&=\mathrm{cmd}.p^{\mathrm{tee}}_{xy}-p^{\mathrm{tee}}_{A,xy},\quad
\Delta_{xy}\leftarrow\Delta_{xy}\cdot\min\!\Big(1,\ \tfrac{\Delta_{\max}}{\|\Delta_{xy}\|_2}\Big),\quad
\Delta_\theta=\mathrm{clamp}\big(\mathrm{wrap}_\pi(\mathrm{cmd}.\theta-\theta_A),\,\pm\theta_{\max}\big),\nonumber\\
\xi&=g(\mathrm{cmd})=\mathrm{clamp}_{\,\mathcal{W}}\Big(p^{\mathrm{tee}}_A+[\Delta_{xy},0],\ \ \theta_A+\Delta_\theta,\ \ q^{\mathrm{full}}_A\Big),
\qquad
(\Delta_{\max},\theta_{\max})=(0.06,0.4).
\label{eq:prescribe}
\end{align}

% (21) The prescribed round: sample the start state from the reset distribution
%      induced by xi, let the expert roll from t*, and aggregate ONE demo.
%      This recovers Diff-DAgger (12) when |S|=1, C_tgt=C*, g=on-policy identity.
\begin{equation}
s_0\sim\Xi(\,\cdot\mid\xi\,),
\qquad
\mathcal{D}_{r+1}=\mathcal{D}_r\ \cup\ \big\{(s_t,\pi^\star(s_t)):\ t\ge t^\star,\ s_0\sim\Xi(\cdot\mid\xi)\big\}.
\label{eq:prescribe-demo}
\end{equation}


% =====================================================================
%  PART V -- EVALUATION METRICS
% =====================================================================

% (22) Held-out success rate over a fixed evaluation set E (target 0.90).
\begin{equation}
\mathrm{SR}(\theta)=\frac{1}{|\mathcal{E}|}\sum_{e\in\mathcal{E}}\mathbf{1}\big[\ \text{episode }e\ \text{succeeds under }\pi_\theta\ \big],
\qquad \mathrm{SR}_{\mathrm{target}}=0.90.
\label{eq:sr}
\end{equation}

% (23) Expert-query count to reach the target SR (else the full budget B).
\begin{equation}
q(\mathrm{SR}_{\mathrm{target}})=\min\big\{\,q\ :\ \mathrm{SR}(\theta_q)\ge\mathrm{SR}_{\mathrm{target}}\,\big\},
\qquad
q\leftarrow B\ \ \text{if the target is never reached.}
\label{eq:queries}
\end{equation}

% (24) Demo efficiency = area under the SR-vs-#demos curve (trapezoidal).
\begin{equation}
\mathrm{Eff}=\int_{0}^{B}\mathrm{SR}(\theta_{q})\ \mathrm{d}q
\ \approx\ \sum_{q=1}^{B}\tfrac12\big(\mathrm{SR}(\theta_{q-1})+\mathrm{SR}(\theta_{q})\big).
\label{eq:eff}
\end{equation}

% (25) Coverage of the collected demos over the task space Z (grid cells in T1;
%      quantized workspace/pose cells in T2-T5).
\begin{equation}
\mathrm{Cov}=\frac{1}{|\mathcal{Z}|}\Big|\ \big\{\,z\in\mathcal{Z}\ :\ \exists\,(s,a)\in\mathcal{D}\ \text{with}\ \mathrm{cell}(s)=z\,\big\}\ \Big|.
\label{eq:cov}
\end{equation}
